% FYE Research & Engagement — Project Proposal
% Template-aligned draft for: Skill Development & Career Readiness for FYE Student Employees
% Note: Several items remain TBD / incomplete pending instrument access, IRB/consent details, and sample sizes.

\documentclass[11pt,letterpaper]{article}

% ---------- Packages (kept minimal for TinyTeX compatibility) ----------
\usepackage[margin=1in]{geometry}
\usepackage{helvet}
\renewcommand{\familydefault}{\sfdefault}
\usepackage[T1]{fontenc}
\usepackage[utf8]{inputenc}
\usepackage{array}
\usepackage{tabularx}
\usepackage{booktabs}

% ---------- Layout helpers ----------
\setlength{\parindent}{0pt}
\setlength{\parskip}{0.4em}
\setlength{\itemsep}{0.15em}
\setlength{\parsep}{0.1em}
\setlength{\topsep}{0.25em}

\newcommand{\sectionrule}{\par\vspace{0.4em}\noindent\rule{\textwidth}{0.6pt}\vspace{0.55em}}
\newcommand{\fieldlabel}[1]{\textbf{#1}}
% Note marker: \tbdtxt{}  or  \tbdtxt{details}
\newcommand{\tbdtxt}[1]{%
  \textit{[TBD\if\relax\detokenize{#1}\relax\else: #1\fi]}%
}

% Template-like numbered section headings
\makeatletter
\renewcommand\section{\@startsection{section}{1}{\z@}%
  {1.0em}{0.4em}{\large\bfseries}}
\makeatother

\begin{document}

% ============================================================
% HEADER
% ============================================================
\begin{center}
  {\LARGE\bfseries FYE Research \& Engagement}\\[0.35em]
  {\Large\bfseries Project Proposal}
\end{center}
\vspace{0.55em}

\fieldlabel{Project Title}\quad Skill Development \& Career Readiness for FYE Student Employees

\vspace{0.45em}
\begin{tabularx}{\textwidth}{@{}l@{\hspace{0.6em}}X@{}}
\fieldlabel{Student Name:} & Pushpak Sunil Rane \\[0.35em]
\fieldlabel{Supervisor:} & Dani Fernandez, FYE Program Coordinator \\[0.35em]
\fieldlabel{Project Duration:} & Aug 24, 2026 \quad to \quad Dec 11, 2026 \\
\end{tabularx}

\sectionrule

% ============================================================
% 1. PURPOSE & OBJECTIVES
% ============================================================
\section{Project Purpose and Objectives}

\fieldlabel{Purpose:}
First-Year Experience (FYE) will assess the skill development and career readiness practices currently built into the program for student employees---including weekly competency reflection, assessment, and feedback; guided professional goal-setting; and interview support. Data will be collected from new FYE student employees, returning employees, and past employees who have started post-graduate careers. Findings will be used to strengthen the program's developmental impact on student employees.

\vspace{0.25em}
\fieldlabel{Objectives:}
By the end of this project, the student will be able to:
\begin{itemize}
  \item Document FYE's existing weekly competency cycle (Practice Plan, Self-Report, supervisor/lead feedback) and related check-in / evaluation touchpoints.
  \item Align FYE competencies with recognized career-readiness frameworks (e.g., NACE Career Readiness Competencies; employer skill priorities from NACE Job Outlook and WEF Future of Jobs sources).
  \item Design and/or refine assessment instruments for weekly reflection quality, mid- and end-of-semester evaluation, goal-setting, and interview readiness \tbdtxt{Work Package fields / competency rubric}.
  \item Collect and analyze multi-cohort evidence (new, returning, and alumni/post-graduate employees) using mixed methods \tbdtxt{survey items, interview protocol, sample size}.
  \item Produce actionable recommendations for FYE to improve competency reflection, feedback, goal-setting, and interview support practices.
\end{itemize}

\sectionrule

% ============================================================
% 2. SCOPE
% ============================================================
\section{Project Scope}

\fieldlabel{In Scope:}
\begin{itemize}
  \item Review of local practice documents (FYE role expectations/Work Package cycle; campus student employment evaluation templates; Career Center handbook resources).
  \item Framework alignment using NACE career readiness competencies / Job Outlook employer attributes and WEF future-skills outlook.
  \item Assessment of weekly competency reflection, written feedback, guided goal-setting, and interview-support activities for FYE student employees.
  \item Multi-cohort inquiry: new FYE student employees, returning employees, and past employees in early post-graduate careers \tbdtxt{sampling / recruitment}.
  \item Synthesis of findings into recommendations for program improvement (instruments, training, feedback quality, interview prep).
\end{itemize}

\vspace{0.2em}
\fieldlabel{Out of Scope:}
\begin{itemize}
  \item Redesign of FYE's full organizational structure, payroll, or hiring policy beyond recommendations related to developmental practices.
  \item Clinical/psychological assessment of students or evaluation of academic course grades outside UNIV~389 / employment-linked reflection artifacts.
  \item Building production software systems or large-scale predictive hiring models.
  \item Generalizable claims beyond FYE / Chico State student-employee context without appropriate caveats.
  \item Experimental interventions that cannot be approved and completed within the semester timeline (deferred to recommendations / future work as needed).
\end{itemize}

\sectionrule

% ============================================================
% 3. METHODOLOGY
% ============================================================
\section{Methodology Overview}

This project uses a \textbf{practice-based, mixed-methods assessment} approach rather than a purely predictive modeling pipeline.

\begin{enumerate}
  \item \textbf{Document \& process mapping.} Inventory FYE's built-in developmental practices (weekly Practice Plan / Self-Report / feedback; 1:1s; mid-semester evaluation; rehire/exit interview; UNIV~389 linkage) and compare with campus Career Center employment evaluation norms.
  \item \textbf{Framework crosswalk.} Map FYE competencies (once confirmed) to NACE Career Readiness Competencies and selected WEF/NACE employer-valued skills (e.g., communication, teamwork, professionalism, critical thinking, leadership, adaptability, lifelong learning; emerging AI literacy as optional).
  \item \textbf{Instrument development / adaptation.} Adapt existing evaluation structures (e.g., 1--4 supervisor rubrics and reflective prompts) and Career Handbook tools (skill statements; STAR interviewing) into study instruments for reflection quality, goal-setting, and interview readiness. \tbdtxt{instrument packet pending Work Package templates}
  \item \textbf{Data collection (planned).}
  \begin{itemize}
    \item Artifact review of de-identified weekly competency entries and feedback samples \tbdtxt{access / consent}
    \item Survey(s) for new and returning employees \tbdtxt{draft items}
    \item Alumni / post-graduate retrospective survey or short interviews on perceived career readiness and use of FYE experiences in interviews \tbdtxt{recruitment}
  \end{itemize}
  \item \textbf{Analysis.} Descriptive quantitative summaries (participation/fidelity, self/supervisor ratings where available) plus qualitative coding of reflection/feedback themes and interview-support usefulness. Triangulate across cohorts where possible.
  \item \textbf{Reporting.} Translate findings into a written report, presentation, and stakeholder recommendations for FYE program improvement.
\end{enumerate}

\vspace{0.3em}
\fieldlabel{Tools may include (but are not limited to):}
\begin{itemize}
  \item Spreadsheet / form tools for rubrics and survey drafts (e.g., Excel, Google Forms, or Qualtrics if available)
  \item Document analysis of FYE Work Package artifacts and evaluation templates
  \item Qualitative coding (manual or tool-assisted spreadsheet codebook)
  \item Descriptive statistics and simple visualizations (Excel, R, or Python)
  \item Reference frameworks: NACE Job Outlook 2026 Spring Update; WEF Future of Jobs Report 2025; Chico State Career Handbook (Vol.~25)
  \item Interview-support materials based on STAR / behavioral interviewing guidance
\end{itemize}

\sectionrule

% ============================================================
% 4. TIMELINE
% ============================================================
\section{Timeline and Benchmarks}

\noindent
\renewcommand{\arraystretch}{1.25}
\begin{tabularx}{\textwidth}{| >{\raggedright\arraybackslash}p{0.34\textwidth} | >{\centering\arraybackslash}p{0.16\textwidth} | >{\raggedright\arraybackslash}X |}
\hline
\textbf{Phase} & \textbf{Timeframe} & \textbf{Key Benchmarks} \\
\hline
Project Planning \& Problem Definition
& Weeks 1--2
& Approved problem statement and project plan; initial document inventory \\
\hline
Practice Mapping \& Framework Review
& Weeks 3--4
& Process map of weekly competency cycle; draft NACE/WEF crosswalk \\
\hline
Instrument Design \& Data Access Prep
& Weeks 5--6
& Draft survey/interview/rubric packet; data-access / consent checklist (\tbdtxt{}) \\
\hline
Data Collection \& Organization
& Weeks 7--10
& Cohort responses and/or artifact sample collected; codebook started (\tbdtxt{depends on access}) \\
\hline
Analysis \& Interpretation
& Weeks 11--12
& Descriptive results + thematic findings; gap analysis vs.\ employer frameworks \\
\hline
Reporting \& Visualization
& Weeks 13--14
& Draft report, figures/tables, and recommendations brief \\
\hline
Final Delivery \& Presentation
& Weeks 15--16
& Final report, presentation, and organized project materials submission \\
\hline
\end{tabularx}
\renewcommand{\arraystretch}{1.0}

\sectionrule

% ============================================================
% 5. DELIVERABLES
% ============================================================
\section{Deliverables}

\fieldlabel{Required Deliverables:}
\begin{enumerate}
  \item \textbf{Project Proposal} (this document)
  \item \textbf{Documented Practice Map \& Competency Crosswalk}
  \begin{itemize}
    \item FYE weekly reflection / feedback / goal-setting / interview-support process map
    \item Alignment matrix to NACE / selected WEF--employer skill priorities
  \end{itemize}
  \item \textbf{Assessment Packet (draft-to-final)}
  \begin{itemize}
    \item Adapted evaluation/reflection instruments
    \item Goal-setting and interview-support prompts/checklists
    \item Notes on incomplete items / assumptions where FYE templates are still pending
  \end{itemize}
  \item \textbf{Dataset(s) / Evidence Binder} (\tbdtxt{scope depends on access})
  \begin{itemize}
    \item Cleaned, documented survey/artifact extracts (de-identified)
    \item Codebook and data dictionary
  \end{itemize}
  \item \textbf{Analysis Workbook / Scripts}
  \begin{itemize}
    \item Well-documented spreadsheets, scripts, or notebooks
    \item Reproducible summary tables/figures where data permit
  \end{itemize}
  \item \textbf{Final Written Report} (15--25 pages, or per program guidelines)
  \begin{itemize}
    \item Executive summary
    \item Methodology
    \item Findings and insights
    \item Limitations and future work
  \end{itemize}
  \item \textbf{Final Presentation}
  \begin{itemize}
    \item 10--15 slide professional presentation summarizing the project
  \end{itemize}
\end{enumerate}

\vspace{0.2em}
\fieldlabel{Optional Deliverables (if applicable):}
\begin{itemize}
  \item Stakeholder-facing recommendations brief for FYE leadership
  \item Interview-support workshop outline for student employees / Training Leads
  \item Interactive summary dashboard (only if dataset size warrants)
  \item Technical appendix (instruments, codebook, additional tables)
\end{itemize}

\sectionrule

% ============================================================
% 6. SUCCESS CRITERIA
% ============================================================
\section{Success Criteria}

The project will be considered successful if:
\begin{itemize}
  \item The assessment approach is appropriate for evaluating FYE's developmental employment practices and is methodologically sound given available data.
  \item FYE practices are clearly documented and meaningfully aligned to external career-readiness / employer skill frameworks.
  \item Results (even if partial) are clearly explained, evidence-supported, and honest about limitations / incomplete elements.
  \item Deliverables meet academic and professional standards expected by FYE Research \& Engagement.
  \item Recommendations are actionable for improving weekly competency reflection, feedback, goal-setting, and/or interview support.
  \item The student demonstrates independent problem-solving, critical thinking, and responsible handling of student-employee information.
\end{itemize}

\sectionrule

% ============================================================
% 7. RISKS
% ============================================================
\section{Risks and Mitigation Strategies}

\noindent
\renewcommand{\arraystretch}{1.3}
\begin{tabularx}{\textwidth}{| >{\raggedright\arraybackslash}p{0.42\textwidth} | >{\raggedright\arraybackslash}X |}
\hline
\textbf{Risk} & \textbf{Mitigation} \\
\hline
Limited access to Work Package artifacts, competency rubrics, or historical evaluations
& Prioritize process documentation + framework crosswalk first; use campus evaluation templates as interim instrument models; escalate data needs early with supervisor \\
\hline
Low survey / alumni response rates
& Keep instruments short; use FYE communication channels; treat qualitative interviews/artifact review as contingency evidence \\
\hline
Scope creep (too many cohorts, instruments, or interventions)
& Advisor check-ins and milestone approvals; freeze MVP instrument set by Weeks 5--6 \\
\hline
Time constraints within the semester
& Clear benchmarks and prioritized tasks; mark incomplete areas explicitly; deliver partial analysis with strong recommendations if data arrive late \\
\hline
Privacy / FERPA / consent constraints
& De-identify data; collect only needed fields; follow supervisor guidance and campus policy before any primary data collection \\
\hline
\end{tabularx}
\renewcommand{\arraystretch}{1.0}

\sectionrule

% ============================================================
% 8. APPROVAL
% ============================================================
\section{Approval}

\vspace{1.2em}
\noindent\textbf{Student Signature:} \underline{\hspace{2.8in}}
\hfill\textbf{Date:} \underline{\hspace{1.3in}}

\vspace{1.5em}
\noindent\textbf{Supervisor Signature:} \underline{\hspace{2.6in}}
\hfill\textbf{Date:} \underline{\hspace{1.3in}}

\vspace{1.4em}
{\footnotesize\textit{Draft proposal aligned to the FYE Research \& Engagement Project Proposal template. Items marked \tbdtxt{} are intentionally incomplete pending further definition with the supervisor (instruments, sampling, data access, and final analytical depth).}}

\end{document}
